\chapter{Computational Study}
\label{sec:comp_study}

In this section, we evaluate the performance of our proposed Itinerary Aggregation Approach (\algoNotation) using an industry-scale instance from Southwest Airlines (WN). We compare the performance of the \algoNotation algorithm against a baseline of solving the original integrated model directly using a state-of-the-art solver (referred to as the ``Raw Model'').

\section{Experimental Setup}
The Southwest Airlines (WN) case study represents a highly connected network characteristic of a major low-cost carrier in the United States. The problem instance is constructed with a 24-hour planning horizon.

The network encompasses $104$ airports and $563$ directed flight segments with a total of $1,119$ scheduled flights per day. The airline operates a homogeneous Boeing 737-based fleet, though $18$ specific fleet sub-types are modeled to account for varying configurations and operational constraints. The seat capacities primarily range from $143$ to $210$ seats (e.g., $143, 175, 210$). 
We utilize a base time discretization of $15$ minutes ($96$ intervals per day) because in the airline industry, 15-minute tolerances are commonly used for scheduling and operational planning.

The passenger demand side includes $1,678$ OD markets. We consider up to $2$ connections with a maximum connection window of $3$ hours. A total of $2,155,200$ feasible itineraries are generated by this process, which are then used as input to the model. This results in a extremely dense and large-scale mathematical program, pushing the limits of conventional memory and computational resources.

\section{Computational Performance}
We compare the convergence behavior and memory footprint of the \algoNotation algorithm against the Raw Model baseline. The results underscore the critical necessity of path-signature-based aggregation and dynamic variable refinement when dealing with millions of itineraries. The performance metrics are summarized in Table~\ref{tab:comp_results}.

\begin{table}[h]
\centering
\caption{Computational Performance: Raw Model vs. \algoNotation (WN Case)}
\label{tab:comp_results}
\begin{tabular}{lrrr}
\hline
\textbf{Method} & \textbf{Final Objective} & \textbf{Optimality Gap} & \textbf{Runtime (s)} \\ \hline
Raw Model (Baseline) & $-1.397 \times 10^7$ & $0.2\%$ & 7,147 \\
\textbf{\algoNotation (Proposed)} & $\mathbf{-1.399 \times 10^7}$ & $< 0.05\%$ & \textbf{2,278} \\ \hline
\end{tabular}
\end{table}

When the complete \problemName model with over 2.15 million itinerary variables is formulated directly and passed to the Gurobi solver, the memory footprint strictly exceeds the available hardware capacity. Consequently, the Raw Model inevitably fails due to an Out of Memory (OOM) error in the early stages of the branch-and-bound process.

In contrast, the \algoNotation algorithm circumvents the memory bottleneck by leveraging the Lower Bound Model (LBM) to solve a dynamically aggregated relaxation. The algorithm then meticulously refines the time-space network and itinerary mappings. 

Within the computational limit, the \algoNotation algorithm successfully discovers a high-quality integer schedule with an objective value of $-13,987,240.63$. Given the solver's predefined relative MIP tolerance of $0.05\%$, the upper and lower bounds perfectly converge, yielding a final optimality gap of \textbf{$< 0.05\%$}. This confirms that the \algoNotation framework mathematically proves and attains the global optimal solution for the entire instance. Resolving an integrated fleet assignment and passenger choice problem of this magnitude to proven optimality underscores the unprecedented scalability of our proposed methodology.

\section{Analysis of the Optimized Schedule}
The globally optimal solution yields profound operational and commercial insights. We systematically extract the activated variables from the exact solution to analyze the underlying scheduling behavior.

The optimal schedule physically activates exactly $1,119$ flights, maintaining active service across all $104$ airports in the network. The fleet assignment strongly favors the 737-8 aircraft, which operates $920$ of the scheduled flights. The remaining $199$ flights are optimally distributed among $13$ different 737-700 series sub-variants (e.g., 737-752, 737-7CT). This heterogeneous deployment ensures that the largest capacities are strictly allocated to routes with sufficient consolidated demand, while older or smaller sub-variants are deployed to efficiently fulfill secondary frequency requirements.

Rather than attempting to blindly serve all potential OD pairs, the algorithm exhibits acute profitability-driven selectivity. Out of the $1,678$ possible markets, the schedule strategically captures demand from only the $496$ most profitable markets, consciously abandoning structurally unprofitable ones to save operating costs. Within these targeted markets, the Sales-Based Linear Program (SBLP) component secures an average market share of $42.6\%$, demonstrating a realistic reflection of endogenous passenger choice and demand spill.

The optimization intricately exploits the network's spatial connectivity without relying on predefined structural rules. The optimal passenger allocation actively utilizes $549$ multi-leg connecting itineraries compared to only $200$ direct itineraries. This heavy reliance on connecting itineraries confirms the model's ability to efficiently stitch together fractional demand from disparate markets. By chronologically synchronizing the departure and arrival bounds of physical flights across the network, the \algoNotation framework seamlessly captures indirect demand and maximizes the total system-wide operating profit.